
\(\underline{O}_X\)-module localement libre de rang \(\underline{n}\), il suffit de montrer, compte tenu de 5.2.2 b), que
\begin{equation}\tag{5.3.3}
 \det(\cl^{\bullet}(\underline{E}))-\cl^{\bullet}(\underline{E})+\underline{n}-1\in\Filt^2(X).
\end{equation}
Comme \(\underline{n}\) est bornée, l'objet final \(\underline{e}\) de \(X\) se décompose en une somme directe \emph{finie} disjointe \(\coprod e_i\) de sous-objets de telle façon que \(\underline{n}\) est constante sur chaque sous-topos \(X/\underline{e}_i\) \((i=1,\ldots,p)\). Comme le foncteur \(Y\rightsquigarrow K^{\bullet}(X/Y)\) \((Y\in\operatorname{Ob}(X))\) transforme sommes directes \emph{finies} en produits directs finis, on peut supposer que \(\underline n\) est constant sur \(X\), soit \(\underline n=n\in\mathbb Z\). Posant \(E=\cl^{\bullet}(\underline E)\), on trouve (cf. Exp V) :
\[
 \det(\underline{E})=\lambda^n(E)=\gamma^n((E-n)+1)
 =
 \sum_{0\leq i\leq n}\gamma^i(E-n),
\]
ce qui démontre (5.3.3), le premier membre étant égal à
\[
 \sum_{2\leq i\leq n}\gamma^i(E-n).
\]

Il nous reste à démontrer le

\textbf{Lemme 5.3.4.} \uline{Pour tout topos annelé} \(X\) \uline{et pour tout} \(x\in\Filt^2(X)\)\uline{, on a}
\[
 \det(x)=1.
\]

\emph{Démonstration} : Soit \(\epsilon:K^{\bullet}(X)\to\Gamma(\mathbb Z)\) l'homomorphisme d'augmentation. Alors (6.2.8) et (6.5.3) ci-dessous donnent par linéarité la formule
\begin{equation}\tag{5.3.5}
 \det(xy)=(\det(x))^{\epsilon(y)}(\det(y))^{\epsilon(x)}
 \qquad(x,y\in K^{\bullet}(X)).
\end{equation}
